\section{Declaration on the candidate’s own contribution}
\label{sec:contributions}
%Purpose of the author’s declaration
%The author’s declaration is submitted to the manuscript committee as an annex to the draft doctoral thesis, to enable the members to assess the PhD candidate’s contribution to the doctoral thesis or to chapters within it. PhD research is a learning process, so the PhD candidate is not expected to do everything completely by himself or herself. For example, the promotor may have conceived the research question, suggested the research approach or suggested improvements to the text. An MSc student may have performed an experiment under the supervision of the PhD candidate. On the basis of the doctoral thesis and the author’s declaration, the manuscript committee will assess whether the doctoral thesis demonstrates the PhD candidate’s competence to conduct independent academic research.
%The author’s declaration focuses on the candidate’s own contribution
%Although it may be necessary to mention other people’s contributions to a chapter, particularly if the PhD candidate is not the first author of a chapter, the declaration should focus on the PhD candidate’s own contribution. The text is therefore to be written in the first person. The text should be concise and should not exceed one A4 page, equivalent to around 500 words. If there are authorship statements for chapters that have already been submitted or published as scientific articles, these can be reused in this author’s declaration. Contributions to the various research chapters typically concern the educational question, the methodology, the research and data collection, the data analysis, the creation of tables and figures, the writing of the text and the provision of feedback on these elements. For other sections or other types of research, it may be relevant to describe other aspects in the author’s declaration.
%The author’s declaration also addresses the use of artificial intelligence tools
%When including a statement on the use of AI in the author’s declaration, it must be made clear whether and for what purpose artificial intelligence tools have been used, for example to improve the text or to generate tables or images. Name the specific AI tools or services used. Indicate that the content has been reviewed and edited, and that full responsibility is assumed for the content of the work. Be transparent about how artificial intelligence tools have been used, for example by listing prompts and applications in footnotes to images. If no artificial intelligence tools have been used, this must also be stated.
Considering the aforementioned research objective, questions, and hypotheses we sum up the research contributions of this thesis below.
\contribution[\acrlong{aar} survey]

\contribution[Sensor-Video Synchronization and Labeling Tool]

To answer \rqname~\ref{rq:labeling}, we discuss a synchronization method to synchronize videos with sensor data.
The method has been used to develop a labeling application that is described in \chaptername~\ref{chapter:data_aquisition_and_labeling}.
Annotations (labels) can be added to the data, and the produced labeled dataset can be exported.
The research and labeling tool have appeared in~\cite{Kamminga_2019_DATA_labeling_sync, Matlab_labeling_github}:

\begin{itemize}
\footnotesize
    \item[\cite{Kamminga_2019_DATA_labeling_sync}] \bibentry{Kamminga_2019_DATA_labeling_sync}
    \item[\cite{Matlab_labeling_github}] \bibentry{Matlab_labeling_github}
\end{itemize}

\subsection{Introduction}
\chaptername~\ref{chapter:state_of_the_art} the research question and the scientific approach were proposed by 
I refined the research question, described how it fills a gap in the academic literature and described its potential societal impact. 
I rewrote the text twice following feedback from my promotor and co-promotors.
\subsection{Chapter2}
\subsection{Chapter3}
    I collected the data, conducted the data analysis and contributed to the revised text of this publication.
    The tables and figures were created by me.
    My co-promotor provided feedback on the interpretation of the results. 
    I rewrote the text XXX-times following feedback from my co-promotor and team.
\subsection{Chapter4}
\subsection{Discussion, conclusion and future work}
    Example: I incorporated the findings from the previous chapters, developed the
    main points of the conclusion and wrote the discussion. I have also made recommendations
    for further research and practical recommendations. The promotor and
    co-promotor provided critical feedback on the content, asked me to refine the
    scientific impact of the doctoral thesis and made suggestions for streamlining the
    text. I rewrote the text twice following feedback from my promotor and co-promotors.
\subsection{Statement on the use of artificial intelligence tools}
    In the preparation of this work I used Grammarly to rewrite the texts and to correct grammatical errors.
    I also used ChatGPT and Visual Studios own coding agend to generate analysis code, bug fix our own analysis tools and adapt this latex code. 
Usage logs, including dates, are available on request.\\[2\baselineskip]

I take full responsibility for the content of this doctoral thesis:
    \vspace{0.5cm}\\
    Enschede, \today\vspace{1cm}\hfill \parbox[t]{5cm}{\centering\hrule\medskip Janike Bolter}

	



